\begin{frame}{Think of the template as code for architecture}
\begin{columns}[T,onlytextwidth]
\column{0.48\textwidth}
\begin{block}{A structured source project}
\begin{itemize}
  \item A main file assembles the manual.
  \item Chapter files separate responsibilities.
  \item Chapter folders hold editable TikZ diagrams.
  \item Build scripts produce the readable PDF.
\end{itemize}
\end{block}
\column{0.48\textwidth}
\begin{block}{A developer's workflow}
\begin{itemize}
  \item Record architecture artifacts in source control alongside the software.
  \item Review text and diagram changes as diffs.
  \item Update design, code, and linked evidence together.
  \item Rebuild and check links before baselining a revision.
\end{itemize}
\end{block}
\end{columns}
\vspace{0.2cm}
\textbf{The file structure organizes the architecture; revisions preserve how it evolves.}
\end{frame}
\note{The analogy is to a software source tree: the entry point includes responsibility-focused modules and shared assets, and a build produces a deliverable. Maintain the editable LaTeX and TikZ sources with the product revision. Record evidence versions and locations so links remain meaningful for that baseline. This is an authoring and configuration-management approach; the architecture itself is not executable software.}
